\documentclass[sigconf, nonacm]{acmart}


\usepackage{draftwatermark}
\SetWatermarkText{AI-Generated}
\SetWatermarkScale{0.5}
\SetWatermarkColor[gray]{0.6}
\AtBeginDocument{\providecommand\BibTeX{{Bib\TeX}}}

\copyrightyear{2025}
\acmYear{2025}
\setcopyright{none}
\acmConference[arXiv preprint]{arXiv preprint}{2025}{}
\acmBooktitle{}
\acmDOI{}
\acmISBN{}
\settopmatter{printacmref=false}
\renewcommand\footnotetextcopyrightpermission[1]{}

\usepackage{multirow}
\usepackage{enumitem}
\usepackage[skip=2pt]{caption}
\usepackage{bm}
\graphicspath{{figures/}}


\begin{document}
\title{FoodLogAthl-218: Constructing a Real-World Food Image Dataset\\ Using Dietary Management Applications}


\author{Anonymous Author}

\begin{abstract}
Food recognition for dietary logging is usually developed on web-crawled images that are attractive, class-balanced, single-dish and independent, whereas photos in a dietary management application are casual, imbalanced, often show several dishes and form a personal timeline. We introduce FoodLogAthl-218, a benchmark built bottom-up from user-submitted meal records of the FoodLog Athl application: labels come from real logs and are cleaned and consolidated afterwards. A pipeline of structural filtering, CLIP-based outlier detection, manual review and GPT-4o-assisted label consolidation turns 11,460 images and 5,889 raw labels from 632 users into 6,925 images, 14,349 dish bounding boxes and 218 classes from 513 users, with timestamps and anonymized user identifiers. The benchmark defines three tasks: closed-set dish classification, continual user-aware adaptation along the time axis, and menu-aware classification of dishes in multi-dish photos. Cluster statistics show larger spread and weaker class separation than UEC-Food 256 in all twelve metric and feature-extractor combinations. In closed-set classification, GPT-4o is strongest (Recall@1 0.424, F1 0.373), no system exceeds 0.60 Recall@5, and the ranking of CLIP-style models and large multimodal models reverses between Recall@1 and Recall@5. Continual LoRA fine-tuning on accumulated logs raises weekly accuracy from roughly 0.26 to roughly 0.49 while zero-shot accuracy falls, a user-specific adapter helps one heavy logger for part of the timeline, and describing the whole meal raises F1 and Precision@1 but lowers Recall@1. We state what the benchmark does not measure, including the absence of inter-annotator agreement and repeated runs.
\end{abstract}

\maketitle
\section{Introduction}
Photo-based dietary logging asks a person to take a picture of a meal and have the system propose what was eaten. The recognition models behind such systems are mostly developed on datasets whose images were collected from the web for a fixed list of categories~\cite{bossard2014food,kawano2015automatic,min2023large}. Web photos are composed to look appealing, categories are chosen first and images collected for them, the class distribution is balanced or curated, each photo usually shows one dish, and images are independent. A photo taken for a private food log differs on every point: it is casual and unfiltered, the label is assigned after the photo exists, class frequency follows what people really eat, one photo often shows a whole meal, and the photos of one person form a timeline with recurring foods. As a result, accuracy on web-sourced benchmarks says little about how a model behaves inside a logging application, and the temporal and personal structure that such applications could exploit~\cite{Horiguchi2018Personalized,Kim2021Boosting} is invisible to them.

We introduce \emph{FoodLogAthl-218}, a benchmark constructed from records of the FoodLog Athl dietary management application~\cite{nakamoto2022foodlog}. The idea is to start from real logs rather than a predefined class list, and then clean and merge labels (Figure~\ref{fig:overview}). The benchmark contains 6,925 images with 14,349 dish bounding boxes in 218 consolidated classes, with anonymized user identifiers, date and time and meal context. It is intended to measure how well a food recognition system copes with the \emph{conditions of real logging}: imbalanced classes, casual imaging, multi-dish meals and temporal continuity. It is not intended to measure nutrition estimation, portion estimation, or recognition outside the cuisine and user population of the source application.

\begin{figure*}[t]
  \centering
  \includegraphics[width=\textwidth]{distribution_5.jpg}
  \caption{Traditional web-sourced food image datasets compared with FoodLogAthl-218 along five properties: label-after-image collection, natural class imbalance, low aesthetic bias, time-series continuity, and multi-dish meals.}
  \label{fig:overview}
\end{figure*}

\paragraph{Design principles.} (i) \emph{Realism}: images and labels originate from users' own records, and the natural long-tailed class distribution is kept. (ii) \emph{Faithful context}: multi-dish photos are kept with all of their dishes (the image is the atomic unit of filtering), and time and user metadata are retained. (iii) \emph{Curated but not sanitized}: automatic outlier detection and manual review remove invalid boxes and labels, and label consolidation reduces synonymous names, but images are not selected for visual appeal. (iv) \emph{Privacy}: two-stage manual removal of personally identifiable information, ethics approval and informed consent.

\paragraph{Findings.} Closed-set classification is hard: GPT-4o, the strongest system, reaches Recall@1 of 0.424 (F1 0.373) and the best Recall@5 is 0.600 (CLIP ViT-L/14). The data are more spread and overlapping than UEC-Food 256 in feature space. Continual LoRA fine-tuning on accumulated logs improves weekly accuracy over a zero-shot model that degrades, and describing the whole meal raises F1 and Precision@1 at the cost of Recall.

\paragraph{Contributions.}
\begin{itemize}[leftmargin=1.2em,itemsep=0pt,topsep=2pt]
  \item A real-world food image dataset built from dietary-app logs, with 6,925 images, 14,349 boxes and 218 classes, natural imbalance and rich metadata.
  \item A documented mixed automatic--human curation pipeline: CLIP outlier filtering, GPT-4o label grouping and manual adjudication.
  \item Three benchmark tasks capturing practical needs: closed-set classification, continual user-aware adaptation, and menu-aware classification of multi-dish meals.
  \item Baselines with CLIP-style models and large multimodal models (LMMs), a feature-space diversity comparison with UEC-Food 256, and an analysis of what adaptation, personalization and meal context change.
\end{itemize}

\section{Related Works}
\paragraph{Food image datasets.} Early datasets were small, such as the Pittsburgh Fast-Food Image Dataset~\cite{chen2009pfid}, Japanese food collections~\cite{joutou2009food,hoashi2010image} and multiple-food recognition~\cite{matsuda2012multiple}. Web crawling made large datasets possible: UEC-Food 256~\cite{kawano2015automatic}, Food-101~\cite{bossard2014food}, UPMC Food-101~\cite{wang2015recipe}, Vegfru~\cite{hou2017vegfru}, FoodAI-756~\cite{sahoo2019foodai}, Food2K~\cite{min2023large} and Food-500 Cap~\cite{ma2023food}, ChineseFoodNet~\cite{chen2017chinesefoodnet}, ISIA Food-500~\cite{min2020isia} and FoodX-251~\cite{kaur2019foodx}. Recipe resources link images to recipes~\cite{salvador2017learning,marin2019learning}, Nutrition5k adds nutritional measurements~\cite{thames2021nutrition5k}, and the Cookpad dataset collects user-posted photos~\cite{harashima2017cookpad}. Table~\ref{tab:datasets} compares the datasets most similar in scope: nearly all draw on the web or on experiment participants, and classes are fixed before images are gathered. FoodLogAthl-218 is modest in size, but among the datasets in the table it is the only one sourced from a dietary management application, where labels are assigned after the photo is taken. See~\cite{DBLP:journals/corr/abs-1808-07202} for a survey.

\paragraph{Dietary logging and personalization.} Photo-based food logging has a long history in multimedia, from FoodLog~\cite{aizawa2015foodlog} and its estimation of food balance from personal dietary tendencies~\cite{Aizawa2013FoodBalance} to FoodLog Athl~\cite{nakamoto2022foodlog}, and mobile systems such as Im2Calories~\cite{meyers2015im2calories} and a recognition engine for dietary logging~\cite{merler2016snap}. Recognition that adapts to the individual has been studied with personalized classifiers~\cite{flg469,Horiguchi2018Personalized} and by sharing food records between users~\cite{Kim2021Boosting}. These works motivate a benchmark that retains user identity and time.

\paragraph{Multimodal models and their evaluation.} Contrastive vision--language models such as CLIP~\cite{CLIP} and EVA-CLIP~\cite{sun2023eva} enable zero-shot classification by image--text matching, and have been examined for Japanese food images~\cite{imajuku2023}. Instruction-tuned LMMs, including the LLaVA family~\cite{liu2024visual,liu2024improved}, Qwen-based models~\cite{bai2023qwen} and GPT-4~\cite{achiam2023gpt}, produce free-form text, which requires a rule to map generations onto a closed label set; we use text-embedding retrieval with E5~\cite{wang2023e5} and the lmms-eval framework~\cite{lmms_eval2024}. Food-specific multimodal assistants~\cite{yin2023foodlmm,Imajuku2025FoodMLLM} and WorldCuisines~\cite{winata2024worldcuisines}, a multilingual visual question answering benchmark on dishes, exist, but none offers real logging statistics, a time axis and per-user structure, which FoodLogAthl-218 is designed to provide.

\section{FoodLogAthl-218 Dataset}
\label{sec:dataset}
\input{tables/table_os}

\subsection{Data source and metadata}
The source is FoodLog Athl, a dietary management application~\cite{nakamoto2022foodlog}. Records span May 2023 to October 2024 and initially comprise 11,460 images with 61,414 bounding boxes and 5,889 raw labels from 632 users. Each record carries an anonymized user ID, date and time, dish names chosen from a database of 140K items, portion size and initial detections; users could correct labels themselves. Because labels are the dish names chosen in the app, the raw label space (5,889 names) is far larger and noisier than a designed class list, and the dataset is therefore constructed bottom-up: images and labels come first, and the class set emerges only after cleaning and consolidation.

\subsection{Curation pipeline}
The image is treated as the atomic unit and multi-dish meals are preserved. Table~\ref{tab:curation} lists the counts after every step. (1--2) Records without an image or without a bounding box are removed, leaving 11,285 images, 26,065 boxes and 3,071 classes. (3) Classes with fewer than 10 samples are dropped (9,121 images, 21,944 boxes, 1,180 classes). (4) CLIP~\cite{CLIP} ViT-B/16 features are used to flag samples further than 1.5 standard deviations from their class centroid, which leaves 7,537 images and 17,513 boxes. (5) Manual review fixes or removes invalid boxes and labels; to maintain the composition of the meal, an image is removed entirely or not at all, giving 6,925 images, 14,349 boxes and 513 users. (6) GPT-4o~\cite{achiam2023gpt} groups synonymous labels, and manual refinement merges them into the final 218 classes (a change of $-962$ classes; the source table printed $-982$, which is inconsistent with its own 1,180 and 218). The review targets incorrect labels, incorrect boxes (e.g., a box covering a different dish) and inappropriate images (e.g., an empty plate labeled ``Milk'' or a product label labeled ``Donut'').

\input{tables/table_zn}

\subsection{Privacy, ethics and quality control}
Personally identifiable information was removed in two manual stages, and the study had ethics approval and informed consent. Quality control combines automatic outlier detection, manual correction of labels and boxes, and human adjudication of the clusters proposed by GPT-4o. Inter-annotator agreement was not measured, so label noise remaining after adjudication cannot be quantified (see Discussion).

\subsection{Resulting dataset}
The retained data cover 513 of the 632 users, i.e., 13.5 images and 28.0 boxes per user on average, and 2.07 boxes per image. Averaged over classes there are 65.8 boxes per class, but the distribution is far from uniform (Figure~\ref{fig:histogram}): no class in the histogram (whose axis starts at 10) has fewer than 10 images, and 16 classes sit at that floor, and the largest class holds roughly $1.7\times10^{3}$ images (read from the plot). Overall, 60.4\% of the images but only 23.4\% of the boxes, and 3.7\% of the raw labels, survive curation. Step 3 alone removes 61.6\% of the classes but only 15.8\% of the boxes that remained after step 2, a direct signature of the long tail of the raw label space.

\begin{figure}[t]
  \centering
  \includegraphics[width=\columnwidth]{histogram_linear_pdf.jpg}
  \caption{Number of dish types (classes) by number of images per class in FoodLogAthl-218 (logarithmic horizontal axis, linear vertical axis). Arrows mark classes at the frequencies shown.}
  \label{fig:histogram}
\end{figure}

\section{Benchmark Tasks and Experimental Setup}
\label{sec:tasks}
\begin{figure}[t]
  \centering
  \includegraphics[width=0.72\columnwidth]{ACM_task.jpg}
  \caption{The three benchmark tasks: (1) closed-set classification of dish crops, where a free-form description generated by the LMM is embedded and matched to class-name embeddings, (2) continual user-aware adaptation along the timeline, (3) menu-aware classification, where a description of the entire image is generated first and then used to predict the dish label for the crop.}
  \label{fig:tasks}
\end{figure}

\subsection{Tasks}
\paragraph{Task 1: closed-set dish classification.} The input is the crop of one annotated dish and the output is one of the 218 classes. CLIP-style models match the crop to class-name text. For LMMs, which generate free text, the model is asked to identify the food, the free-form answer is embedded with \texttt{intfloat/e5-large-v2}~\cite{wang2023e5}, and the class whose label embedding has the highest cosine similarity is chosen; ranked similarities give Top-$K$ predictions. LMMs are run through lmms-eval~\cite{lmms_eval2024}.

\paragraph{Task 2: continual user-aware adaptation.} Because records carry time stamps, data arrive in order. A model is updated on records accumulated up to a point in time and evaluated on subsequent records (Figure~\ref{fig:tasks}, bottom left). Accuracy is reported for each week, together with a weighted one-month moving average over $\pm2$ weeks, whose weights are proportional to the weekly data volume. Two settings are used: adaptation on the logs of all users, and adaptation on the logs of a single user, which can be compared with the all-user model on that user's later records.

\paragraph{Task 3: menu-aware classification.} Meals often show several dishes, and what else is on the tray is informative about the dish in a crop. In this task the LMM first produces a description of the whole image (Figure~\ref{fig:description}) and the crop is then classified given that description; the setting is compared with classifying the crop alone.

\begin{figure}[t]
  \centering
  \includegraphics[width=0.85\columnwidth]{description_eg.jpg}
  \caption{Description-guided classification. Left: the crop alone is misclassified (Okonomiyaki for Dumplings) and the description of the whole image fixes it. Right: the crop alone is classified correctly (Vegetable\_Juice), but the description leads to a wrong answer (Egg Soup).}
  \label{fig:description}
\end{figure}

\subsection{Metrics and systems}
We report Top-$K$ accuracy as Recall@$K$ (R@1, R@3, R@5), Precision@1 (P@1) for Task 3, and F1 computed from top-1 predictions. Recall@$K$ rewards a correct class anywhere in the top $K$, which suits a logging interface that shows several candidates, while F1 and P@1 reflect how often the single top prediction is right across classes. The systems are GPT-4o~\cite{achiam2023gpt}; open LMMs from the Qwen2-VL and Qwen2.5-VL families, LLaVA-OneVision, LLaVA-v1.5 and LLaVA-v1.6~\cite{liu2024visual,liu2024improved,bai2023qwen}; and the contrastive models CLIP ViT-L/14~\cite{CLIP} and EVA-CLIP ViT-g~\cite{sun2023eva}. Task 2 uses LoRA~\cite{hu2022lora} fine-tuning of Qwen2.5-VL-3B and full fine-tuning for comparison, and Task 3 uses Qwen2.5-VL-7B-Instruct. All numbers below come from single evaluation runs; no confidence intervals or repeated seeds are available, and we do not add any.

\section{Benchmark Validation}
\label{sec:validation}
\input{tables/table_xj}

\paragraph{Distribution compared with a web-sourced dataset.} To check that the real-log data are harder than a web-sourced benchmark in feature space, we compare against UEC-Food 256~\cite{kawano2015automatic} using CLIP~\cite{CLIP}, EVA-CLIP~\cite{sun2023eva} and Swin~\cite{liu2021swin} features and four statistics: Trace, Silhouette~\cite{rousseeuw1987silhouettes}, Davies--Bouldin~\cite{davies1979cluster} and Calinski--Harabasz~\cite{calinski1974dendrite} (Table~\ref{tab:diversity}). FoodLogAthl-218 has larger spread and weaker class separation in all twelve extractor--metric combinations: for CLIP, Trace 81.70 versus 71.21, Silhouette $-0.021$ versus $0.034$, Davies--Bouldin 4.52 versus 3.64 and Calinski--Harabasz 27.51 versus 74.05. Negative Silhouette values with CLIP and Swin mean that, on average, samples are closer to another class than to their own. Two caveats apply: the datasets differ in class number (218 versus 256), size and content, and these statistics are descriptive, not a calibrated measure of task difficulty.

\paragraph{Difficulty and non-triviality.} The closed-set results (next section) show that no system exceeds a Recall@5 of 0.600 or an F1 of 0.373, so the benchmark is far from saturated, and that systems are separated: R@1 ranges from 0.193 to 0.424 over ten models.

\paragraph{Label and annotation reliability.} Evidence for label quality is procedural (outlier filtering, manual review with all-or-none image removal, human adjudication of the GPT-4o synonym clusters), not statistical: no inter-annotator agreement or audit-sample error rate is available. Merging 1,180 classes into 218 reduces ambiguity from synonymous names but cannot remove inherent ambiguity between visually similar dishes; the failure in Figure~\ref{fig:description} (right) is an example of a plausible confusion between a drink and a soup.

\section{Results and Analysis}
\label{sec:results}
\subsection{Task 1: closed-set classification}
\input{tables/table_kf}
Table~\ref{tab:classification} reports Task 1. GPT-4o is best on F1 (0.373), R@1 (0.424) and R@3 (0.499), and CLIP ViT-L/14 is best on R@5 (0.600), followed by EVA-CLIP (0.594). The weakest systems, llava-v1.5-7b and llava-v1.6-mistral-7b, reach R@1 of 0.214 and 0.193. Three patterns stand out.

\emph{The metric used changes the ranking.} The CLIP-style models are 6th and 7th of ten by R@1 (0.265 and 0.257) but 1st and 2nd by R@5, and their R@5 is 0.329 (EVA-CLIP) and 0.343 (CLIP) above R@1, whereas for the LMMs the gap is only 0.083 to 0.163. Over the ten systems the Spearman rank correlation is 0.53 between R@1 and R@5 and 0.94 between R@5 and F1 (computed from Table~\ref{tab:classification}). Which behavior is preferable depends on whether an application shows one suggestion or a list.

\emph{F1 and R@1 diverge for LMMs.} For GPT-4o, F1 is 0.88 of R@1, but for Qwen2-VL-7B it is 0.44, for llava-onevision-qwen2-7b-ov 0.43 and for llava-v1.5-7b 0.46. Qwen2.5-VL-7B is considerably higher in F1 (0.235) than Qwen2-VL-7B (0.159) with almost the same R@1 (0.368 versus 0.366). We did not analyze per-class errors; a possible explanation, correct predictions concentrated in a few frequent classes, is a hypothesis, not a result.

\emph{Scale alone does not explain the order.} llava-v1.6-34b (R@1 0.243) is below the 7B llava-onevision-qwen2-7b-ov (0.363), and Qwen2-VL-3B (0.266) is below Qwen2-VL-7B (0.366) in the same family.

\subsection{Task 2: continual user-aware adaptation}

Figure~\ref{fig:allusers} shows weekly accuracy and its moving average from May 2023 to January 2024 (values read from the plot, hence approximate). The zero-shot moving average starts near 0.23 and declines to about 0.16, while LoRA fine-tuning on accumulated logs rises from about 0.26 to about 0.49 and full fine-tuning from about 0.24 to about 0.51. The LoRA model thus ends about 33 percentage points above zero-shot and about 23 above its own start; full fine-tuning is ahead for most of the period but the two are close at the end. Raw weekly points fluctuate by 0.1 or more (small weekly volumes), so differences of a few points between variants should not be over-interpreted. The decline of zero-shot accuracy, which we cannot explain with the available data, shows that the benchmark exposes change along the timeline, which a static test set cannot.

\begin{figure*}[t]
  \centering
  \begin{minipage}[t]{0.49\textwidth}\centering
    \includegraphics[width=\linewidth]{all_data_1_35_with_legend.jpg}
    \caption{Task 2, all users: weekly accuracy (dashed) and weighted moving average (solid) for zero-shot inference, full fine-tuning and LoRA fine-tuning on accumulated logs.}
    \label{fig:allusers}
  \end{minipage}\hfill
  \begin{minipage}[t]{0.49\textwidth}\centering
    \includegraphics[width=\linewidth]{user_zero_lora_ma_wide_legend.jpg}
    \caption{Task 2, one heavy logger (User A): zero-shot, LoRA fine-tuning on all users' data and LoRA fine-tuning on User A's data only.}
    \label{fig:usera}
  \end{minipage}
\end{figure*}
Figure~\ref{fig:usera} shows one heavy logger from May to December 2023 (approximate readings). Zero-shot accuracy stays between roughly 0.05 and 0.14. The all-user LoRA model rises from about 0.30 to about 0.65. The user-specific adapter starts lower (about 0.17), is ahead of the all-user model from about August to November 2023 (e.g., about 0.52 versus 0.39 around September), peaks near 0.67, and ends below it (about 0.57 versus 0.65). Relative to zero-shot at the end, the user-specific adapter gains roughly 50 percentage points and the all-user adapter about 55. For this one user, personalization thus helps for part of the timeline but does \emph{not} end above adaptation on all users' data; with one user and strongly fluctuating weekly points we draw no general conclusion about personalization.

\subsection{Task 3: menu-aware classification}
\input{tables/table_vs}
Table~\ref{tab:desc_guided} shows that providing the description of the whole meal raises F1 from 0.159 to 0.230 (+0.071) and P@1 from 0.102 to 0.179 (+0.077), while R@1, R@3 and R@5 fall from 0.366 to 0.324, from 0.459 to 0.403 and from 0.478 to 0.434. The context makes the top-1 prediction more precise but loses correct answers overall. At the level of individual predictions, the description corrects 687 errors and introduces 1,222 new ones (a net 535 additional errors), so the context helps some dishes, as in Figure~\ref{fig:description} (left), and hurts others (right). Menu-level context is therefore not a free improvement, and the benchmark makes the trade-off measurable. One caution on the data: the ``w/o description'' row of Table~\ref{tab:desc_guided} (F1 0.1594, R@1 0.3656, R@3 0.4590, R@5 0.4781) is numerically identical to the Qwen2-VL-7B-Instruct row of Table~\ref{tab:classification}, whereas the Qwen2.5-VL-7B-Instruct row there differs (F1 0.2349, R@1 0.3684). We report both tables as provided and could not resolve which of them carries a mislabeled model name; the direction of the effect of the description does not depend on this, but the identity of the baseline row does.

\section{Discussion}
\paragraph{What the benchmark offers.} FoodLogAthl-218 tests food recognition under real logging conditions through three views (static classification, continual and personal adaptation, meal context), and its results are informative beyond a single score: system ordering depends on R@1, R@5 or F1, and meal context shifts the precision--recall balance.

\paragraph{What it does and does not support.} Conclusions on food logging in the population and cuisine of one application are supported; conclusions on other cuisines, on nutrition estimation, or on user-facing quality are not. The benchmark is not comprehensive for ``food recognition'', and a high score on it would not imply that a system is safe to use for dietary advice.

\paragraph{Limitations and risks.} (i) Inter-annotator agreement was not measured and the label set results from GPT-4o grouping plus manual refinement, so merged classes reflect judgment calls. (ii) Curation discards three quarters of the raw boxes; the survivors are cleaner than the logs, and the outlier filter, which uses CLIP features, may favor images that CLIP finds typical, a possible bias toward CLIP-style models that we have not measured. (iii) The data come from a single application and user population, with 13.5 images per user on average, and single users dominate the personalization analysis. (iv) The supplied results include no confidence intervals, seeds, or repeated runs, so the Task 2 and Task 3 values are descriptive. (v) The comparison to UEC-Food 256 is not controlled for class number or content. (vi) Splits, the GPT-4o version, prompts, the averaging of F1 and fine-tuning hyperparameters are not part of the material available to this write-up, so full reproduction needs the original experiment code and logs, which were not supplied. (vii) Foundation models may have seen similar food photos; contamination was not tested. (viii) A row shared between the Task 1 and Task 3 tables is inconsistent (Section~\ref{sec:results}).

\section{Conclusion}
We presented FoodLogAthl-218, a food image benchmark built from dietary-app logs with 6,925 images, 14,349 dish boxes and 218 classes, and three tasks that reflect real logging: closed-set classification, continual user-aware adaptation and menu-aware classification. The data are more spread and overlapping than UEC-Food 256 in all twelve cluster-metric comparisons, closed-set recognition is far from solved (best R@1 0.424, best R@5 0.600) and system ranking depends on the metric. Continual fine-tuning improves weekly accuracy over a decaying zero-shot baseline, user-specific adaptation helps one heavy logger for part of the timeline, and meal-level descriptions trade recall for precision. Given missing agreement statistics, a single data source and single-run evaluations, the benchmark enables research on continual, personalized and context-aware food recognition under realistic conditions, not claims of general food understanding.

\bibliographystyle{ACM-Reference-Format}

\bibliography{references}

\end{document}
